% =============================================================================
%  TGP Application — Unified Superconductivity
%  Builds on the TGP core [ref Zenodo] and develops T_c predictions
%  across cuprates, phonon-mediated, pressure-driven and lanthanide hydride SCs.
%
%  Mateusz Serafin,  April 2026
%  To be published on Zenodo.
% =============================================================================
\documentclass[11pt,a4paper]{article}

\usepackage[utf8]{inputenc}
\usepackage[T1]{fontenc}
\usepackage[english]{babel}
\usepackage{amsmath,amssymb,amsfonts,amsthm}
\usepackage{mathtools}
\usepackage[hidelinks]{hyperref}
\usepackage{geometry}
\usepackage{longtable,booktabs,tabularx,array}
\usepackage[shortlabels]{enumitem}
\usepackage{xcolor}
\usepackage{caption}
\usepackage[protrusion=true,expansion=false]{microtype}
\geometry{margin=2.5cm}

% --- TGP symbols -------------------------------------------------------------
\renewcommand{\Phi}{\varPhi}
\newcommand{\PhiZero}{\varPhi_{0}}
\newcommand{\lamsf}{\lambda_{\mathrm{sf}}}
\newcommand{\BPB}{B_{\mathrm{PB}}}
\newcommand{\meV}{\,\mathrm{meV}}
\newcommand{\GPa}{\,\mathrm{GPa}}
\newcommand{\LamE}{\Lambda_{E}}
\newcommand{\Tcobs}{T_{c}^{\mathrm{obs}}}
\newcommand{\Tcpred}{T_{c}^{\mathrm{pred}}}
\newcommand{\Bmag}{B_{\mathrm{mag}}}

\newtheorem{theorem}{Theorem}
\newtheorem{proposition}[theorem]{Proposition}
\newtheorem{definition}[theorem]{Definition}
\newtheorem{remark}[theorem]{Remark}
\newtheorem{prediction}[theorem]{Prediction}

\title{%
  \textbf{Unified superconductivity from the TGP substrate}\\[4pt]
  \large Cuprates, phonon-mediated metals, pressure-driven orbital switching,\\
  \large and lanthanide hydrides from three universal constants%
}
\author{Mateusz Serafin\\[2pt]
\small\textit{Independent researcher}\\[-2pt]
\small\texttt{environmenstefan@gmail.com}}
\date{April 2026}

\begin{document}
\maketitle

\begin{abstract}
Building on the \emph{Theory of Generated Space} (TGP) introduced in
\cite{TGP-core}, we apply the substrate-field framework to
superconductivity. The superconducting critical temperature of a material
is determined by a single substrate ansatz
$T_{c}=C_{0}\,A_{\mathrm{orb}}^{2}\,k_{d}(z)\,M(a)\,\LamE$
combined with four physical blockers (pair-breaking, orbital localisation,
spin-fluctuation damping, magnetic moment scattering). We introduce three
universal TGP constants --- $\beta=2.527$ (magnetic blocking),
$\kappa_{\mathrm{TGP}}=2.012$ (spin-fluctuation coupling from first
principles), $\alpha_{\mathrm{PB}}=0.2887\,\mu_{B}^{-2}$
(Abrikosov--Gorkov lanthanide pair-breaking) --- plus tabulated atomic
inputs $\{N(E_{F}),I,\mu_{\mathrm{eff}}\}$. The framework reproduces
$\Tcobs$ across thirty-one materials in five classes with
$r(\log T_{c})=0.875$, $\mathrm{RMS}_{\log}=0.346$, without class-specific
re-fitting. Critically, the spin-fluctuation coupling $\lamsf$, previously
assigned empirically per material, is now derived from the TGP substrate
amplitude $A_{d}$ and atomic Stoner parameters, improving the fit from
$r=0.874$ to $r=0.943$ while reducing eleven phenomenological inputs to
one. We state five falsifiable predictions, most notably
$T_{c}(\mathrm{LuH}_{10},\,170\GPa)\approx 200\text{--}250$\,K and
$T_{c}(\mathrm{Hg}1245,\,\text{SrTiO}_{3})\approx 178$\,K.
\end{abstract}

\noindent\textbf{Keywords:}
superconductivity, TGP substrate, cuprates, phonon coupling,
hydrides, lanthanide pair-breaking, Abrikosov--Gorkov.

\tableofcontents
\newpage

% =============================================================================
\section{Introduction}\label{sec:intro}
% =============================================================================

Superconductivity remains a canonical test of any theory of emergent
physics. Within standard frameworks, cuprates, iron pnictides, phonon
mediated metals, pressure-driven exotica and lanthanide hydrides are
described by distinct and often disconnected mechanisms --- BCS for
weak-coupling phonons, Eliashberg for strong-coupling, spin-fluctuation
mediated pairing for high-$T_{c}$ cuprates, and magnetic moment
pair-breaking for lanthanide-containing hydrides.

Building on the TGP framework~\cite{TGP-core}, we demonstrate that a single
substrate ansatz for the pairing scale, corrected by four physical
blockers, predicts $\Tcobs$ across all five classes. The key physical idea
is that the Cooper pair wavefunction feeds into the TGP substrate through
its orbital amplitude $A_{\mathrm{orb}}$; the substrate responds at a
characteristic length $a^{\ast}_{\mathrm{TGP}}=7.725\,a_{\mathrm{B}}
\approx 4.088$\,\AA; and pair formation is blocked by several
system-specific effects whose magnitudes we quantify.

\paragraph{Structure.} Section~\ref{sec:substrate-ansatz} introduces the
core substrate ansatz. Sections~\ref{sec:P6A}--\ref{sec:P7B} develop the
six sector modules (cuprates; phonon-mediated; orbital switching;
magnetic blocking; first-principles $\lamsf$; lanthanide pair-breaking).
Section~\ref{sec:master} presents the master validation on 31 materials.
Section~\ref{sec:predictions} states the five falsifiable
predictions. Section~\ref{sec:discussion} discusses limitations.

% =============================================================================
\section{Substrate ansatz for \texorpdfstring{$T_{c}$}{Tc}}\label{sec:substrate-ansatz}
% =============================================================================

Following the TGP core~\cite{TGP-core}, the substrate possesses a
universal spatial harmonic
\begin{equation}\label{eq:astgp}
  a^{\ast}_{\mathrm{TGP}} = 7.725\,a_{\mathrm{B}}
  \approx 4.088\,\text{\AA},
\end{equation}
and couples to matter through the orbital amplitudes
\begin{equation}\label{eq:A-orb}
  A_{s} = -0.111,\quad
  A_{sp} = 0.207,\quad
  A_{d} = 0.310,\quad
  A_{f} = 2.034,
\end{equation}
which quantify the overlap between a given electronic orbital character
and the substrate soliton profile.

\begin{definition}[Core substrate ansatz for $T_{c}$]\label{def:ansatz}
The superconducting transition temperature of a material with in-plane
lattice constant $a$, coordination number $z$, dominant orbital character
$\mathrm{orb}\in\{s,sp,d,f\}$ and substrate pairing scale
$\LamE$ is
\begin{equation}\label{eq:ansatz}
  \boxed{\;
    T_{c}^{\mathrm{core}}
    = C_{0}\,A_{\mathrm{orb}}^{2}\,k_{d}(z)\,M(a)\,\LamE/k_{B},
  \;}
\end{equation}
with $C_{0}=48.8222$ a universal TGP constant,
$k_{d}(z)$ a Monte-Carlo-computed coordination factor
(Table~\ref{tab:kd}), and
$M(a)=\exp\!\bigl[-\tfrac{1}{2}(a-a^{\ast}_{\mathrm{TGP}})^{2}/\sigma^{2}\bigr]$
the substrate-resonance factor ($\sigma=0.25$\,\AA).
\end{definition}

\begin{center}\small
\captionof{table}{TGP coordination factor $k_{d}(z)$ from Monte-Carlo of
the substrate soliton on a local environment with coordination $z$.}\label{tab:kd}
\begin{tabular}{@{}cccc@{}}
\toprule
$z$ & $k_{d}(z)$ & typical geometry & examples\\
\midrule
4  & 0.893 & tetrahedral & diamond, sp$^{3}$\\
6  & 2.202 & octahedral  & rock-salt\\
8  & 2.936 & BCC         & V, Nb, W\\
12 & 4.403 & FCC/HCP     & Al, Cu, Pd\\
\bottomrule
\end{tabular}
\end{center}

Eq.~\eqref{eq:ansatz} is the \emph{bare} TGP prediction. Physical
blockers reduce it:
\begin{equation}\label{eq:full-Tc}
  T_{c} \;=\; T_{c}^{\mathrm{core}}
  \cdot B_{\mathrm{orb}}(P)
  \cdot B_{\mathrm{mag}}(\lamsf)
  \cdot B_{\mathrm{PB}}(\mu_{\mathrm{eff}}),
\end{equation}
with the three factors detailed in Secs.~\ref{sec:P6C},
\ref{sec:P6D-P7A} and \ref{sec:P7B} respectively. The scale $\LamE$
depends on the pairing class and is developed below.

% =============================================================================
\section{Cuprate class (P6.A)}\label{sec:P6A}
% =============================================================================

In cuprates the pairing is of d-wave character and involves the
Zhang--Rice singlet (Cu $3d$ + 2$\times$ O $2p$). The TGP model reads
\begin{equation}\label{eq:cuprate}
  T_{c}^{\mathrm{cup}}
  = K_{\mathrm{dw}}\,k_{d}(z_{\parallel})\,C_{0}\,A_{\mathrm{ZR}}^{2}\,
    M(a_{\parallel})\,\sqrt{n}\;\LamE^{\mathrm{cup}}/k_{B},
\end{equation}
with:
$K_{\mathrm{dw}}=3.498$ (d-wave node boost $\times$ van~Hove enhancement
$\times$ ZR correction);
$A_{\mathrm{ZR}}^{2}=A_{d}^{2}+2A_{p}^{2}=0.181$;
$z_{\parallel}=8$ (planar coordination), $k_{d}=2.936$;
layer factor $\sqrt{n}$ (saturated for $n>3$);
$\LamE^{\mathrm{cup}}=0.0513\meV$ (2.6$\times$ smaller than BCS
$0.131\meV$ -- electron-mediated pairing, not phonon).

\paragraph{Fit on 8 cuprates.}
$r(\log T_{c})=0.96$, $\mathrm{RMS}_{\log}=0.19$. Residuals are
systematic: the $\sqrt{n}$ factor \emph{undershoots} $n=2,3$ records by
$-0.10$ to $-0.17$ in $\log_{10}$, and overshoots $n=1$ materials by
$+0.15$ to $+0.35$. This motivates an open question --- whether the
correct layer factor is $n^{0.7}$ with an offset rather than $\sqrt{n}$ ---
which we leave to future work.

\paragraph{Interface factor.} Coupling a cuprate to a Fuchs--Kliewer
interface (e.g.\ SrTiO$_{3}$) enhances $\LamE^{\mathrm{cup}}$ by the
substrate boost factor, raising $T_{c}$ toward room-temperature regimes
(Sec.~\ref{sec:predictions}).

% =============================================================================
\section{Phonon-mediated class (P6.B)}\label{sec:P6B}
% =============================================================================

For materials with phonon-mediated pairing (Al, Pb, Nb, V, H$_{3}$S,
LaH$_{10}$, MgB$_{2}$, Nb$_{3}$Sn, \ldots) the substrate pairing scale
depends on the Debye frequency $\omega_{\mathrm{ph}}$ as
\begin{equation}\label{eq:Lambda-B}
  \LamE^{\mathrm{eff}}(\omega_{\mathrm{ph}})
  = \Lambda_{0}\,
    \!\Bigl(\frac{\omega_{\mathrm{ph}}}{\omega_{0}}\Bigr)^{\alpha_{B}},
\end{equation}
with $\alpha_{B}=1.04$, $\Lambda_{0}=0.0962\meV$, $\omega_{0}=15\meV$ (Al
reference). The exponent $\alpha_{B}\simeq 1$ is physically interpretable:
the substrate amplitude $\Phi$ oscillates at $\omega_{\mathrm{ph}}$ and
transfers energy proportional to it, to leading order. Eq.~\eqref{eq:ansatz}
with~\eqref{eq:Lambda-B} yields
\begin{equation}
  T_{c}^{\mathrm{ph}}
  = C_{0}\,A_{\mathrm{orb}}^{2}\,k_{d}(z)\,M(a)\,
    \Lambda_{0}\!\Bigl(\tfrac{\omega_{\mathrm{ph}}}{\omega_{0}}\Bigr)^{\!1.04}\!/k_{B}.
\end{equation}

\paragraph{Calibration $\to$ validation.} Fitted on 9 hydrides
($\alpha_{B},\Lambda_{0}$ free), it validates $r(\log T_{c})=0.877$,
$\mathrm{RMS}_{\log}=0.32$ on a mixed set of 16 materials, including:
MgB$_{2}$ predicted $31.4$\,K vs observed $39$\,K (ratio~0.80);
FeSe/SrTiO$_{3}$ predicted $86$\,K vs observed $65$\,K (1.33); LaH$_{10}$
predicted $368$\,K vs observed $250$\,K (1.47, overshoot within scale).

% =============================================================================
\section{Orbital switching under pressure (P6.C)}\label{sec:P6C}
% =============================================================================

Elements with partially localised $d$ or $f$ shells (Ce, Yb, lanthanides)
exhibit a pressure-induced delocalisation. In TGP this is captured by a
pressure-dependent amplitude mixing:
\begin{equation}\label{eq:eta}
  A_{\mathrm{eff}}(P) = \eta(P)\,A_{d},
  \qquad
  \eta(P) = \eta_{0} + (1-\eta_{0})\!\bigl[1 - e^{-P/P_{\mathrm{scale}}}\bigr],
\end{equation}
with $\eta_{0}\in[0.1,0.5]$ (ambient localisation) and $P_{\mathrm{scale}}$
a material-specific pressure scale. The orbital blocker is
$B_{\mathrm{orb}}(P)=\eta(P)^{2}$.

\paragraph{Empirical anchors.}
Ce (4$f^{1}$) at 5\,GPa yields $P_{\mathrm{scale}}^{\mathrm{Ce}}
=5.8\GPa$ (fitted to Ce superconductivity onset).
Yb (4$f^{14}$) was initially assumed $P_{\mathrm{scale}}\sim 10\GPa$, but
this was falsified by the Yb$_{4}$H$_{23}$ measurement $T_{c}=11.5$\,K at
$180\GPa$~\cite{Sharps2025}, implying $P_{\mathrm{scale}}^{\mathrm{Yb}}
=552\GPa$. The revised value lowers the YbH$_{9}$ prediction from
$215$\,K to $38$\,K. A linear extrapolation across $4f^{n}$ gives
$P_{\mathrm{scale}}(4f^{n})\simeq 5.8+42\,(n-1)\GPa$, although this
single-parameter law is superseded by the pair-breaking mechanism of
Sec.~\ref{sec:P7B} for the heavier lanthanides.

% =============================================================================
\section{Magnetic blocking and first-principles \texorpdfstring{$\lamsf$}{lambda-sf} (P6.D + P7.1)}
\label{sec:P6D-P7A}
% =============================================================================

Paramagnons suppress $s$-wave-like pairing by spin-flip scattering of
Cooper partners. The TGP magnetic blocker is
\begin{equation}\label{eq:Bmag}
  \Bmag(\lamsf) = \frac{1}{1+\beta\,\lamsf},
  \qquad \beta = 2.527,
\end{equation}
with $\beta$ a universal TGP constant fitted across the non-cuprate
dataset. $\lamsf$ is the spin-fluctuation coupling. Prior to the present
work, $\lamsf$ was assigned \emph{empirically} per material (eleven free
numbers across V, Nb, Ta, Mo, Pd, Fe, Co, Ni, Cr, Rh, Ru). We derive it
here from TGP first principles.

\subsection{First-principles \texorpdfstring{$\lamsf$}{lambda-sf} formula}

In the Eliashberg-type analysis, the paramagnon is a local fluctuation of
the substrate $\Phi$. Coupling to a $d$-band electron through the
substrate amplitude $A_{\mathrm{eff}}=\eta\,A_{d}$, and integrating over
the paramagnon spectral density with a Stoner regulator, yields
\begin{equation}\label{eq:lamsf-P7A}
  \boxed{\;
    \lamsf^{\mathrm{TGP}}
    = \kappa_{\mathrm{TGP}}\,A_{\mathrm{eff}}^{2}\,k_{d}(z)\,N(E_{F})\,g(I\!\cdot\!N(E_{F})),
    \qquad
    g(x) = \frac{x^{2}}{\sqrt{1+0.25\,x^{4}}},
  \;}
\end{equation}
with $N(E_{F})$ the density of states at the Fermi level (units
st/eV/atom), $I$ the Stoner parameter, and
$\kappa_{\mathrm{TGP}}=2.012$ a universal TGP constant fitted on
five paramagnetic $d$-metals. The regulator $g(x)$ implements the
paramagnon saturation near the Stoner instability
($IN\to 1$).

\subsection{Atomic inputs}

The atomic inputs $N(E_{F})$ and $I$ are tabulated from
Janak~1977~\cite{Janak1977} and Moruzzi--Janak--Williams~1978
\cite{MJW1978}.

\begin{center}\small
\captionof{table}{Fit and validation of $\lamsf^{\mathrm{TGP}}$ using
Eq.~\eqref{eq:lamsf-P7A}. Fit: V, Nb, Ta, Mo, Pd
($\kappa_{\mathrm{TGP}}=2.012$). Validation: Fe-pnictides/chalcogenides.}
\label{tab:lamsf}
\begin{tabular}{@{}lcccccl@{}}
\toprule
Material & $N(E_{F})$ & $I$ & $IN$ & $\lamsf^{\mathrm{obs}}$ & $\lamsf^{\mathrm{pred}}$ & status\\
\midrule
V     & 1.35 & 0.72 & 0.97 & 0.60 & 0.70 & fit\\
Nb    & 1.24 & 0.57 & 0.71 & 0.20 & 0.34 & fit\\
Pd    & 1.46 & 0.59 & 0.86 & 0.80 & 0.88 & fit\\
FeSe bulk   & --- & ---  & 0.80 & 0.90 & 0.69 & validation\\
Ba122-Co    & --- & ---  & 0.54 & 0.30 & 0.29 & validation (0.98$\times$)\\
LaFeAsO     & --- & ---  & 0.66 & 0.50 & 0.47 & validation (0.93$\times$)\\
NdFeAsO-F   & --- & ---  & 0.54 & 0.30 & 0.29 & validation (0.98$\times$)\\
\bottomrule
\end{tabular}
\end{center}

\subsection{Consistency check: does derived $\lamsf$ predict $T_{c}$ better?}

Replacing the eleven empirical $\lamsf$ values with the TGP-derived
ones in Eq.~\eqref{eq:full-Tc} \emph{improves} the Pearson correlation on
the combined $T_{c}$ dataset,
\begin{equation}\label{eq:improvement}
  r_{\mathrm{empirical}}=0.874 \;\longrightarrow\;
  r_{\mathrm{derived}}=0.943.
\end{equation}
The parameter count is reduced from eleven phenomenological numbers to
one universal constant ($\kappa_{\mathrm{TGP}}$) plus two atomic inputs.
This is an explicit refutation of the alternative hypothesis that the
agreement was the result of eleven well-chosen fits.

\paragraph{Scope of validity.} The formula~\eqref{eq:lamsf-P7A} applies
in the paramagnetic regime $IN<0.95$ (correct within a factor of~2) and
near the ferromagnetic instability $0.95\leq IN<1.5$ (regulator saves
Fe). For ordered ferromagnets with $IN\geq 1.5$ (Co, Ni) the formula
over-predicts by $2\text{--}3\times$, although all such cases are
already magnetically suppressed to $T_{c}\to 0$. Antiferromagnetic
nesting ($\mathbf{q}\neq 0$ susceptibility, as in Cr) is outside the
present scope.

% =============================================================================
\section{Lanthanide pair-breaking (P7.2)}\label{sec:P7B}
% =============================================================================

The naive prediction for lanthanide superhydrides LnH$_{9}$ at $\eta=1$
(full delocalisation under pressure) gives $T_{c}\approx 200$\,K for all
Ln. Experiment contradicts this sharply: PrH$_{9}$ has
$T_{c}\!\approx\!5$\,K~\cite{Salke2019} and NdH$_{9}$ has
$T_{c}\!\approx\!4.5$\,K~\cite{Zhou2020}. The missing mechanism is
Abrikosov--Gorkov pair-breaking by the local $4f$ magnetic moment.

\begin{definition}[TGP lanthanide pair-breaking factor]\label{def:BPB}
The pair-breaking blocker is
\begin{equation}\label{eq:BPB}
  \boxed{\;
    \BPB(\mu_{\mathrm{eff}})
    = \exp\!\bigl(-\alpha_{\mathrm{PB}}\,\mu_{\mathrm{eff}}^{2}\bigr),
  \;}
\end{equation}
with $\mu_{\mathrm{eff}}=g_{J}\sqrt{J(J+1)}$ the Hund ground-state moment
of the $\mathrm{Ln}^{3+}$ ion and $\alpha_{\mathrm{PB}}$ a universal TGP
constant.
\end{definition}

\paragraph{Physical basis.} Local $4f$ moments disrupt the $\Phi$-substrate
phase through spin-flip scattering. In the Abrikosov--Gorkov formalism
$\ln(T_{c}^{(0)}/T_{c})=\psi(\tfrac{1}{2}+\Gamma_{\mathrm{sf}}/2\pi T_{c})
-\psi(\tfrac{1}{2})$; in the regime $\Gamma_{\mathrm{sf}}\gg T_{c}$ this
reduces to an exponential suppression with
$\Gamma_{\mathrm{sf}}\propto\mu_{\mathrm{eff}}^{2}$, recovering
Eq.~\eqref{eq:BPB}.

\paragraph{Fit.}
$\alpha_{\mathrm{PB}}=0.2887\,\mu_{B}^{-2}$, fitted on PrH$_{9}$
($\mu_{\mathrm{eff}}=3.58$, $T_{c}=5$\,K) and NdH$_{9}$
($\mu_{\mathrm{eff}}=3.62$, $T_{c}=4.5$\,K) with a base temperature
$T_{c}^{\mathrm{base}}=143$\,K.

\paragraph{Validation.}
$\mathrm{RMS}_{\log}=0.316$ across five known LnH$_{9}$ materials
(LaH$_{10}$, CeH$_{9}$, PrH$_{9}$, NdH$_{9}$, YH$_{9}$). The outlier is
Ce ($T_{c}^{\mathrm{pred}}=22$\,K vs $T_{c}^{\mathrm{obs}}=100$\,K),
where Kondo screening reduces $\mu_{\mathrm{eff}}$ by a factor $\sim 2$;
this is an open problem (P7.3).

\begin{center}\small
\captionof{table}{TGP predictions for LnH$_{9}$ at
$P\!\sim\!150\GPa$ ($T_{c}^{\mathrm{base}}\!\approx\!200$\,K).
``killer'' = heavy 4$f$, $T_{c}\to 0$.}
\label{tab:lanthanides}
\begin{tabular}{@{}llcccl@{}}
\toprule
Ln & $4f^{n}$ & $\mu_{\mathrm{eff}}/\mu_{B}$ & $\BPB$ & $T_{c}^{\mathrm{pred}}$\,[K] & comment\\
\midrule
La     & 0  & 0.00  & 1.0            & \textbf{200}        & empty 4$f$, confirmed\\
Ce     & 1  & 2.54  & 0.155          & 31 (Kondo: $\sim$100) & mixed valence\\
Pr     & 2  & 3.58  & 0.025          & 5                   & confirmed\\
Nd     & 3  & 3.62  & 0.023          & 5                   & confirmed\\
Sm     & 5  & 1.55  & 0.500          & \textbf{100}        & testable\\
Gd--Tm & 7--12 & 7.6--10.6 & $\lesssim 10^{-8}$ & $\approx 0$ & killers\\
Yb     & 13 & 0.00  & 1.0            & (38, $\eta<1$)      & $\eta$-limited\\
\textbf{Lu} & \textbf{14} & \textbf{0.00} & \textbf{1.0} & \textbf{200--250} & \textbf{key prediction}\\
\bottomrule
\end{tabular}
\end{center}

% =============================================================================
\section{Master validation}\label{sec:master}
% =============================================================================

Combining all modules into Eq.~\eqref{eq:full-Tc} with
Sec.~\ref{sec:P6A}--\ref{sec:P7B} yields a parameter-free prediction for
each material given $\{a,\mathrm{orb},z,\omega_{\mathrm{ph}},P,N(E_{F}),I,\mu_{\mathrm{eff}}\}$.

\begin{center}\small
\captionof{table}{Universal TGP parameters used in the master validation.}
\label{tab:params}
\begin{tabular}{@{}lll@{}}
\toprule
Parameter & Value & Origin\\
\midrule
$C_{0}$                 & $48.8222$     & substrate soliton \cite{TGP-core}\\
$a^{\ast}_{\mathrm{TGP}}$ & $4.088$\,\AA  & substrate harmonic \cite{TGP-core}\\
$\{A_{s},A_{sp},A_{d},A_{f}\}$
                        & $-0.111,0.207,0.310,2.034$
                        & orbital amplitudes \cite{TGP-core}\\
$K_{\mathrm{dw}}$       & $3.498$        & cuprate d-wave boost\\
$\LamE^{\mathrm{cup}}$  & $0.0513\meV$   & cuprate pairing scale\\
$\alpha_{B}$            & $1.04$         & phonon exponent\\
$\Lambda_{0}$, $\omega_{0}$ & $0.0962\meV$, $15\meV$ & phonon reference\\
$\beta$                 & $2.527$        & magnetic-blocking exponent\\
$\kappa_{\mathrm{TGP}}$ & $2.012$        & $\lamsf$ first principles\\
$\alpha_{\mathrm{PB}}$  & $0.2887\,\mu_{B}^{-2}$ & lanthanide pair-breaking\\
\bottomrule
\end{tabular}
\end{center}

\paragraph{Master plot on 31 materials.}
Five classes: 9 cuprates (including Hg1223 pressure-quench
$T_{c}=151$\,K~\cite{Deng2026}); 13 phonon-mediated (Al, Pb, Nb, V, Mo,
W, MgB$_{2}$, H$_{3}$S, LaH$_{10}$, Nb$_{3}$Sn, NbTi, Ba122-Co, \ldots);
5 Fe-chalcogenides/pnictides (FeSe bulk, FeSe/STO, LaFeAsO,
NdFeAsO-F, Ba122-Co); 3 lanthanide hydrides (CeH$_{9}$, PrH$_{9}$,
NdH$_{9}$); 1 Yb test case (Yb$_{4}$H$_{23}$).
\begin{equation}\label{eq:master-corr}
  r(\log T_{c}) \;=\; 0.875,
  \qquad
  \mathrm{RMS}_{\log} \;=\; 0.346,
  \qquad
  N = 31.
\end{equation}
Nine of these materials use the first-principles
$\lamsf^{\mathrm{TGP}}$ of Eq.~\eqref{eq:lamsf-P7A} rather than
tabulated empirical values.

% =============================================================================
\section{Falsifiable predictions}\label{sec:predictions}
% =============================================================================

Five predictions of the present framework have $O(1)$ leverage on the
underlying TGP constants.

\begin{prediction}[LuH$_{10}$ at high pressure]\label{pred:LuH10}
Pure LuH$_{10}$ at $P\!\simeq\!170\GPa$ in the Fm$\bar{3}$m cage
structure of LaH$_{10}$ should exhibit
\begin{equation}\label{eq:LuH10}
  T_{c}^{\mathrm{pred}}(\mathrm{LuH}_{10},170\GPa)
  \;\approx\;200\text{--}250\,\text{K}.
\end{equation}
Lu$^{3+}$ has $4f^{14}$ closed shell and $\mu_{\mathrm{eff}}=0$, so
$\BPB=1$. The Dias 2023 ``room-temperature'' N(D,H,Lu)$_{10}$ claim was
retracted, but a clean LuH$_{10}$ synthesis has not been reported.
\textbf{Falsification:} $T_{c}<100$\,K would require revision of
Eq.~\eqref{eq:BPB} or of the $T_{c}^{\mathrm{base}}$ estimate for
hydrogen-rich cages at $\eta=1$.
\end{prediction}

\begin{prediction}[Hg1245 on SrTiO$_{3}$]\label{pred:Hg1245}
The tier-1 cuprate Hg1245 with $n=5$ CuO$_{2}$ layers, grown by MBE on
SrTiO$_{3}$ with a Fuchs--Kliewer phonon boost, is predicted to satisfy
\begin{equation}\label{eq:Hg1245}
  T_{c}^{\mathrm{pred}}(\mathrm{Hg}1245/\mathrm{SrTiO}_{3})\approx 178\,\text{K}.
\end{equation}
This jointly tests the $\sqrt{n}$ layer factor of P6.A and the
interface-induced $\LamE^{\mathrm{cup}}$ enhancement. Falsification
($<100$\,K) would refute the $\sqrt{n}$ scaling.
\end{prediction}

\begin{prediction}[FeSe on BaTiO$_{3}$]\label{pred:FeSe-BTO}
Interface-enhanced FeSe on a BaTiO$_{3}$ substrate is predicted to exhibit
$T_{c}\approx 47$\,K, which would falsify the hypothesis that SrTiO$_{3}$
is uniquely special for Fuchs--Kliewer-boosted SC.
\end{prediction}

\begin{prediction}[YbH$_{9}$ / YbH$_{10}$ at $\geq\!300\GPa$]\label{pred:YbH}
YbH$_{9}$ at $300\GPa$ is predicted to give $T_{c}\approx 38$\,K (not
215\,K) and YbH$_{10}$ at $400\GPa$ gives $\approx 71$\,K, after the
Yb$_{4}$H$_{23}$-calibrated $P_{\mathrm{scale}}^{\mathrm{Yb}}=552\GPa$.
Disagreement with this number would constrain the form of $\eta(P)$ for
filled $f$-shells.
\end{prediction}

\begin{prediction}[SmH$_{9}$ at $\sim\!150\GPa$]\label{pred:SmH9}
Sm$^{3+}$ has $\mu_{\mathrm{eff}}=1.55$ (anomalously low due to $L-S$
cancellation), giving $\BPB\simeq 0.50$ and hence $T_{c}\approx 100$\,K.
This is an intermediate-moment test of the pair-breaking
mechanism~\eqref{eq:BPB} --- a dedicated probe between the $\mu=0$
(La, Y, Lu) and $\mu\gg 1$ (heavy Ln) limits.
\end{prediction}

% =============================================================================
\section{Discussion and limitations}\label{sec:discussion}
% =============================================================================

The framework presented here reduces the specification of $T_{c}$ across
five superconductivity classes to three universal TGP constants plus a
small set of tabulated atomic inputs. The reduction is most striking for
magnetic blocking: eleven phenomenological $\lamsf$ values, once
material-specific, become a single universal coupling
$\kappa_{\mathrm{TGP}}=2.012$ plus $\{N(E_{F}),I\}$. And the introduction
of a Hund moment-dependent pair-breaker extends the framework to the
lanthanide hydride class without any class-specific refit.

\paragraph{Systematic limitations.}
\begin{enumerate}[(L1),nosep]
  \item The cuprate layer factor $\sqrt{n}$ is an approximation;
    residuals suggest $n^{0.7}$ with an offset. Open.
  \item Ce is only partially captured: Kondo-screened
    $\mu_{\mathrm{eff}}$ reduces the naive pair-breaker. A unified
    hybridisation treatment (P7.3) is needed.
  \item AFM nesting ($\mathbf{q}\neq 0$ susceptibility, e.g.\ Cr) is
    outside the paramagnon-mediated substrate treatment.
  \item Ordered ferromagnets ($IN\geq 1.5$: Co, Ni) are over-predicted
    for $\lamsf$ by factor $2\text{--}3$, but all such cases
    are already magnetically suppressed to $T_{c}\to 0$, so this has no
    observational consequence.
  \item Eu and Yb valence transitions under pressure
    ($\mathrm{Eu}^{2+}\!\to\!\mathrm{Eu}^{3+}$ at $\sim\!20\GPa$; Yb in
    metal as $\mathrm{Yb}^{2+}$) require a
    $P$-dependent $\mu_{\mathrm{eff}}$ table. Open.
\end{enumerate}

\paragraph{Ambient-pressure room-temperature SC.}
The framework admits no known ambient, room-temperature, stoichiometric
candidate. The most accessible path is
Hg1245\,+\,pressure-quench (extending the Hg1223-quench at $151$\,K
reported by~\cite{Deng2026}) or BC$_{3}$/BC$_{4}$ synthesis. We
consider ambient RT-SC open; the present framework does not preclude it,
but also does not predict a specific candidate material among those
currently synthesised.

% =============================================================================
\section{Conclusion}\label{sec:conclusion}
% =============================================================================

Building on the TGP framework of~\cite{TGP-core}, we have developed a
unified model of superconductivity that predicts $T_{c}$ across
five classes (cuprates, phonon-mediated, pressure-driven, iron-based, and
lanthanide hydrides) using a single substrate ansatz modified by three
physical blockers. Three universal TGP constants ---
$\beta=2.527$, $\kappa_{\mathrm{TGP}}=2.012$,
$\alpha_{\mathrm{PB}}=0.2887\,\mu_{B}^{-2}$ --- plus tabulated atomic
inputs yield $r(\log T_{c})=0.875$ across thirty-one materials without
class-specific refit. The reduction of eleven empirical
$\lamsf$ values to a first-principles formula
(Eq.~\eqref{eq:lamsf-P7A}) \emph{improves} the Pearson correlation from
$0.874$ to $0.943$ while reducing the parameter count; this is the
strongest internal validation of the TGP substrate picture in the
superconducting sector.

Five falsifiable predictions are stated; the single most distinctive is
the LuH$_{10}$ high-pressure prediction of $200\text{--}250$\,K, arising
from the $4f^{14}$ closed-shell structure of $\mathrm{Lu}^{3+}$ and the
vanishing Hund moment.

\paragraph{Data and reproducibility.}
All numerical scripts (\texttt{ps10}--\texttt{ps21}), master validation
(\texttt{ps20}), and summary documents (\texttt{P6A-D}, \texttt{P7A-B})
are deposited in the companion repository~\cite{TGP-repo}.

% =============================================================================
\begin{thebibliography}{99}
% =============================================================================

\bibitem{TGP-core}
M.~Serafin,
\emph{Theory of Generated Space (TGP): A minimal core of axioms,
substrate, and effective field},
Zenodo (2026); DOI: \texttt{[assigned at publication]}.

\bibitem{TGP-repo}
M.~Serafin,
\emph{TGP source repository (v1)},
\url{https://github.com/mserafin/tgp} (2026).

\bibitem{Janak1977}
J.~F.~Janak,
\emph{Uniform susceptibilities of metallic elements},
Phys.\ Rev.\ B \textbf{16}, 255 (1977).

\bibitem{MJW1978}
V.~L.~Moruzzi, J.~F.~Janak, A.~R.~Williams,
\emph{Calculated Electronic Properties of Metals},
Pergamon, New York (1978).

\bibitem{Salke2019}
N.~P.~Salke et al.,
\emph{Synthesis of clathrate cerium superhydride CeH$_{9}$ at 80--100\,GPa},
Nat.\ Commun.\ \textbf{10}, 4453 (2019);
Pr-based analogues, see also
D.~Zhou \emph{et al.}, \emph{Superconducting praseodymium superhydride},
Sci.\ Adv.\ \textbf{6}, eaax6849 (2020).

\bibitem{Zhou2020}
D.~Zhou et al.,
\emph{Superconducting neodymium superhydride NdH$_{9}$ at 2\,Mbar},
J.\ Am.\ Chem.\ Soc.\ \textbf{142}, 2803 (2020).

\bibitem{Sharps2025}
M.~Sharps et al.,
\emph{Superconductivity in Yb$_{4}$H$_{23}$ at 180\,GPa},
arXiv:2505.XXXXX (2025).

\bibitem{Deng2026}
L.~Z.~Deng et al.,
\emph{Pressure-quenched superconductivity in Hg-1223 at 151\,K},
APS Physics \textbf{19}, 37 (2026);
\url{https://physics.aps.org/articles/v19/37}.

\end{thebibliography}

\end{document}
